\subsection*{§. 197.\\ 运动作为绝对形到自身的射影变换。}

椭圆空间的每个运动都是表示空间的一个射影变换，并把绝对二次曲面变为自身；因而，若考察绝对形的两族母线，则每一族都射影地变为自身\srcfn{1}{连续运动群显然不能交换两族母线。}。特别在一个滑移中，由于所有点都沿直线运动，其中一族的各条母线必须分别保持不变；不妨设为第一族。第二族的母线则彼此作射影交换，并有两条特殊的（共轭）母线保持不动，显然连同其每一点一起保持不动。因此，一个滑移无非是以同一族的两条共轭母线为轴的双轴射影变换；反过来，立即可见每个这样的双轴射影变换都是一个滑移。

因此，一个 Clifford 直线汇的 $\infty^2$ 条直线，其特征就是都与绝对形的两条共轭母线相交。

容易注意到，两个同类滑移复合后仍是同类滑移；而变换一个滑移 $S$ 时，若以任意运动 $M$ 来作这一变换，即构成运动 $M^{-1}SM$，仍得到同类滑移 $S'$\srcfn{2}{事实上，运动 $M^{-1}SM$ 保持滑移 $S$ 所保持的同一族母线不变，所以是同类滑移。}。

由此可知，同类的 $\infty^3$ 个滑移在六参数总运动群中组成一个三参数不变子群。

再者，容易看出，任何运动都能分解为两个异类滑移的相继作用。事实上，在不是滑移的任意运动 $M$ 中，第一族只有两条（共轭）母线 $g_1,g_2$ 保持不动，另一族也只有两条 $g'_1,g'_2$ 保持不动。这四条母线组成一个空间四边形，其两条对角线是关于二次曲面互为极线的实直线，并在运动中分别沿自身滑动；亦即每个运动都有两条关于绝对形互为极线的运动轴 $a,a'$。现在取第一族中另一条母线 $g_3$（不同于 $g_1,g_2$），它在 $M$ 下移到 $g'_3$；把 $M$ 与第二类滑移 $S'$ 复合，该滑移保持 $g'_1,g'_2$ 不动并把 $g'_3$ 移到 $g_3$。复合运动 $MS'$ 保持三条母线 $g_1,g_2,g_3$ 不动，因而保持这一族的所有其他母线不动，所以是第一类滑移 $S$；于是
\[
 MS'=S,
 \qquad M=SS'^{-1}.
\]
还要注意，两个滑移 $S,S'$，一个为第一类、一个为第二类，总是彼此可交换的，即
\[
 SS'=S'S,
\]
因为运动 $SS'$ 与运动 $S'S$ 对同一族母线产生相同作用，因而二者重合。此外，也容易用解析方法验证这里以几何方法导出的运动性质。
